\section*{Capítulo \ref{ch:g11:mutations} --- Mutaciones y variación genética}

\begin{solution}{exo:g11:mutations:1}
Un cambio en la secuencia de \omterm{def:g10:universal-dna:nucleotide}{nucleótidos} del \omterm{def:g10:universal-dna:information}{ADN}, transmitido a los
descendientes de la \omterm{def:g10:cells-common-unit:cell}{célula}. Sustitución, inserción, deleción.
\end{solution}

\begin{solution}{exo:g11:mutations:2}
\texttt{ATGCCAGAAT}: sustitución (T por A en la posición 6).
\texttt{ATGCCTGAT}: deleción de una A. \texttt{ATGCCTTGAAT}: inserción
de una T.
\end{solution}

\begin{solution}{exo:g11:mutations:3}
Somática: en una \omterm{def:g10:cells-common-unit:cell}{célula} del cuerpo, transmitida solo a los
descendientes de esa \omterm{def:g10:cells-common-unit:cell}{célula} dentro del individuo. Germinal: en un
gameto o en su precursor, transmitida a la descendencia y a todas sus
\omterm{def:g10:cells-common-unit:cell}{células}. Solo las \omterm{def:g11:mutations:mutation}{mutaciones} germinales entran en la herencia.
\end{solution}

\begin{solution}{exo:g11:mutations:4}
La luz ultravioleta, los rayos X y el benzopireno del humo del tabaco
(también la radiactividad y la aflatoxina). El ultravioleta suelda
entre sí dos timinas vecinas de una hebra y deforma la hélice.
\end{solution}

\begin{solution}{exo:g11:mutations:5}
La hebra complementaria intacta: como cada hebra determina la otra, la
secuencia correcta del segmento dañado puede reconstruirse a partir de
su pareja.
\end{solution}

\begin{solution}{exo:g11:mutations:6}
Normal: alrededor del 6\%; deficiente en reparación: 0,01\%. La
reparación mejora la supervivencia unas 600 veces a esa dosis.
\end{solution}

\begin{solution}{exo:g11:mutations:7}
$20/10^9 = \num{2e-8}$. No: la placa solo revela qué \omterm{def:g10:cells-common-unit:cell}{células} eran ya
resistentes; para saber si el antibiótico causó la \omterm{def:g11:mutations:mutation}{mutación} hace falta
la prueba de la réplica.
\end{solution}

\begin{solution}{exo:g11:mutations:8}
Cada réplica recibe \omterm{def:g10:cells-common-unit:cell}{células} de las mismas colonias originales. Si el
antibiótico causara la resistencia al contacto, las \omterm{def:g10:cells-common-unit:cell}{células} afectadas
serían unas pocas al azar en cada réplica, en posiciones sin relación
entre sí. Encontrar la resistencia siempre en las mismas posiciones
significa que aquellas colonias originales, crecidas sin el fármaco, ya
contenían \omterm{def:g10:cells-common-unit:cell}{células} resistentes.
\end{solution}

\begin{solution}{exo:g11:mutations:9}
Errores por copia del \omterm{def:g10:universal-dna:gene}{gen}: $2 \times 3000 \times 10^{-9} =
\num{6e-6}$ (las dos hebras). Entre los $10^{12}$ \omterm{def:g10:cells-common-unit:cell}{células} producidas en
la última división, unas $\num{6e6}$ llevan una \omterm{def:g11:mutations:mutation}{mutación} nueva en ese
\omterm{def:g10:universal-dna:gene}{gen}.
\end{solution}

\begin{solution}{exo:g11:mutations:10}
Una \omterm{def:g11:mutations:mutation}{mutación} somática en una \omterm{def:g10:cells-common-unit:cell}{célula} de yema dio su color a la rama; los
esquejes son ese tejido y por eso lo conservan. Las \omterm{def:g10:cells-common-unit:cell}{células} germinales
de las flores proceden de esa misma rama, pero un cambio en un \omterm{def:g10:universal-dna:gene}{alelo}
solo pasa a la mitad de los gametos, y los óvulos fueron fecundados por
polen con el \omterm{def:g10:universal-dna:gene}{alelo} corriente; el \omterm{def:g10:universal-dna:gene}{alelo} nuevo, si es recesivo, queda
oculto en las plántulas.
\end{solution}

\begin{solution}{exo:g11:mutations:11}
La \omterm{def:g10:cell-metabolism:metabolism}{enzima} que falta repara el daño ultravioleta, y el ultravioleta solo
penetra en la piel. El tubo digestivo no recibe ultravioleta, no sufre
esas lesiones y no necesita esa reparación.
\end{solution}

\begin{solution}{exo:g11:mutations:12}
Las lesiones son proporcionales a la dosis: una décima parte del
ultravioleta significa una décima parte de las lesiones por \omterm{def:g10:cells-common-unit:cell}{célula} y
por hora --- que siguen siendo miles. La quemadura solar es la muerte
de muchas \omterm{def:g10:cells-common-unit:cell}{células}; por debajo de esa dosis las \omterm{def:g10:cells-common-unit:cell}{células} sobreviven, pero
las lesiones que escapan a la reparación se convierten en \omterm{def:g11:mutations:mutation}{mutaciones}.
Que no haya quemadura significa que el daño era soportable, no que no
lo hubiera.
\end{solution}

\begin{solution}{exo:g11:mutations:13}
Ninguna de las dos cosas por sí sola. En una región con paludismo, una
sola copia aumenta la supervivencia, y el \omterm{def:g10:universal-dna:gene}{alelo} se mantiene frecuente
pese a la enfermedad de quienes llevan dos; donde no hay paludismo solo
causa enfermedad y es raro. El ambiente, y el número de copias, deciden
el valor de una \omterm{def:g11:mutations:mutation}{mutación}.
\end{solution}

\begin{solution}{exo:g11:mutations:14}
Las \omterm{def:g10:cells-common-unit:cell}{células} sensibles mueren; las resistentes, ya presentes, se dividen
sin trabas y en un día el cultivo entero es resistente. Una tanda
interrumpida antes de tiempo deja vivos a los últimos supervivientes
--- los más resistentes --- para que vuelvan a crecer en un cuerpo ya
enriquecido en resistencia; completar la tanda los mata antes de que se
multipliquen.
\end{solution}

\begin{solution}{exo:g11:mutations:15}
Casi todas las \omterm{def:g11:mutations:mutation}{mutaciones} son neutras (fuera de los \omterm{def:g10:universal-dna:gene}{genes} o sin efecto
sobre la \omterm{def:g10:chemistry-of-life:families}{proteína}), una minoría es perjudicial y muy pocas son útiles;
al ser aleatorias, ocurren con independencia de su utilidad, pero la
selección conserva las útiles y elimina las perjudiciales. La utilidad
depende del ambiente, así que una \omterm{def:g11:mutations:mutation}{mutación} perjudicial hoy puede ser la
que importe cuando cambien las condiciones. A lo largo de miles de
generaciones, con miles de millones de individuos, incluso las
\omterm{def:g11:mutations:mutation}{mutaciones} útiles raras surgen una y otra vez: la evolución se
construye sobre el residuo seleccionado, no sobre la \omterm{def:g11:mutations:mutation}{mutación} media.
\end{solution}

\begin{solution}{pb:g11:mutations:1}
\textbf{1.} Para que cada colonia sea un clon de origen conocido, y
para que las colonias estén lo bastante separadas como para
distinguirse en las réplicas.

\textbf{2.} Unas pocas \omterm{def:g10:cells-common-unit:cell}{células} de la parte superior de cada colonia, en
la misma disposición geométrica que en la placa, porque la almohadilla
se aprieta en plano y sin deslizarse.

\textbf{3.} Que las \omterm{def:g10:cells-common-unit:cell}{células} resistentes estaban presentes en aquellas
dos colonias originales: tres réplicas independientes no coinciden por
azar.

\textbf{4.} Colonias resistentes en posiciones aleatorias y sin
relación entre sí en las tres placas, ya que qué \omterm{def:g10:cells-common-unit:cell}{células} quedaran
inducidas variaría de una réplica a otra.

\textbf{5.} \omterm{def:g11:mutations:mutation}{Mutación} antes del contacto: la resistencia se localiza en
las colonias originales, que nunca se encontraron con el fármaco.

\textbf{6.} Las \omterm{def:g10:cells-common-unit:cell}{células} de las dos posiciones resistentes crecen con
estreptomicina; las de las otras posiciones, no.

\textbf{7.} Probablemente solo algunas: la \omterm{def:g11:mutations:mutation}{mutación} ocurrió en alguna
división durante el crecimiento de la colonia, y solo los
descendientes de aquella \omterm{def:g10:cells-common-unit:cell}{célula} son resistentes. Para decidirlo habría
que sembrar un número conocido de \omterm{def:g10:cells-common-unit:cell}{células} de la colonia con
estreptomicina y contar los supervivientes.

\textbf{8.} En la división 20 de 23: $1/2^{20-1}$\dots más
sencillamente, la \omterm{def:g10:cells-common-unit:cell}{célula} mutante se divide 3 veces más: $2^3 = 8$
\omterm{def:g10:cells-common-unit:cell}{células} resistentes de $2^{23}$, alrededor de una entre un millón. En
la división 3: una de las 8 \omterm{def:g10:cells-common-unit:cell}{células} presentes entonces, o sea una
octava parte de la colonia.

\textbf{9.} Bastan unas pocas \omterm{def:g10:cells-common-unit:cell}{células} resistentes entre las $10^4$
aproximadamente que levanta el terciopelo para fundar una colonia en la
placa con antibiótico; la prueba detecta la presencia de \omterm{def:g10:cells-common-unit:cell}{células}
resistentes, no su proporción.

\textbf{10.} El argumento descansa en la posición: solo si se conserva
la geometría puede identificarse una colonia de una réplica con una
colonia de la placa original.

\textbf{11.} $30/\num{10000} = \num{3e-3}$.

\textbf{12.} Unas $\num{3e-3}$ \omterm{def:g11:mutations:mutation}{mutaciones} por cada $10^7$ divisiones:
$\num{3e-10}$ por división.

\textbf{13.} Del mismo orden que la tasa de error en un \omterm{def:g10:universal-dna:nucleotide}{nucleótido}: la
\omterm{def:g11:mutations:mutation}{mutación} de resistencia es esencialmente una sustitución concreta que
escapa a la reparación.

\textbf{14.} $10^9$ divisiones a $\num{3e-10}$: de media, una fracción
de \omterm{def:g11:mutations:mutation}{mutación}, así que casi todos los cultivos no contienen ninguna
\omterm{def:g10:cells-common-unit:cell}{célula} resistente o contienen unas pocas (cada \omterm{def:g11:mutations:mutation}{mutación} multiplicada
después por las divisiones que la siguen).

\textbf{15.} Una \omterm{def:g11:mutations:mutation}{mutación} en una división temprana la heredan una gran
fracción del cultivo, así que el recuento es alto; una tardía da solo
unas pocas \omterm{def:g10:cells-common-unit:cell}{células} resistentes. Como el momento es aleatorio, los
cultivos repetidos dan recuentos muy dispersos --- lo cual es en sí
mismo un indicio de que las \omterm{def:g11:mutations:mutation}{mutaciones} surgen antes del antibiótico y
no como respuesta a él.

\textbf{16.} Raras: las \omterm{def:g11:mutations:mutation}{mutaciones} ocurren pocas veces por división.
Aleatorias: ocurren con independencia de su utilidad. El experimento
establece la aleatoriedad (el argumento de las posiciones); los
recuentos de la Parte III miden la rareza.

\textbf{17.} Clasificó la población: mató a las \omterm{def:g10:cells-common-unit:cell}{células} sensibles y
dejó multiplicarse a las resistentes --- selección.

\textbf{18.} La \omterm{def:g11:mutations:mutation}{mutación} aportó el \omterm{def:g10:universal-dna:gene}{alelo} de resistencia, al azar, antes
de toda necesidad; el antibiótico clasificó la población y conservó
solo a quienes lo llevaban.

\textbf{19.} La \omterm{def:g11:mutations:mutation}{mutación} cambia una molécula implicada en la acción del
fármaco, no el crecimiento ni el aspecto de la colonia sin él; una
diferencia genética solo tiene efecto visible en un ambiente en el que
importa.

\textbf{20.} Las posiciones idénticas en las tres réplicas demostraron
que las \omterm{def:g11:mutations:mutation}{mutaciones} de resistencia surgen antes de la exposición al
antibiótico y con independencia de ella; la tasa medida es de unos
$\num{3e-10}$ por división celular.
\end{solution}
